\chapter{Projeto do Sistema}
\label{cap:projeto}

\orientacao{Este capítulo define como a solução será estruturada para atender aos requisitos identificados na análise. Sua função é transformar o modelo conceitual em uma especificação técnica da solução, estabelecendo sua arquitetura, componentes, dados, interfaces, interações, tecnologias e mecanismos de segurança.}

\section{Arquitetura do Sistema}
\orientacao{Apresentar a organização geral da solução, definindo seus principais elementos, responsabilidades e relações em nível arquitetural.}

\diagrama{Diagrama de componentes}{componentes}

\section{Projeto dos Componentes}
\orientacao{Detalhar a estrutura e as responsabilidades dos principais componentes definidos na arquitetura, incluindo suas interfaces e relações.}

\section{Modelo de Dados}
\orientacao{Descrever a estrutura dos dados necessários ao funcionamento do sistema e suas relações, considerando as decisões de projeto adotadas.}

\diagrama{Diagrama de classes de projeto}{classes-projeto}

\section{Projeto da Interface do Usuário}
\orientacao{Apresentar a organização e o comportamento das interfaces pelas quais os usuários interagem com o sistema, incluindo protótipos ou representações visuais quando pertinentes.}

\section{Diagramas de Interação}
\orientacao{Detalhar como os elementos do sistema colaboram para realizar os comportamentos definidos, evidenciando a sequência de mensagens e responsabilidades.}

\diagrama{Diagrama de sequência}{sequencia}

\section{Tecnologias Utilizadas}
\orientacao{Relacionar as principais linguagens, frameworks, bibliotecas, sistemas gerenciadores, ferramentas e demais tecnologias adotadas, indicando sua função no projeto.}

A \autoref{tab:tecnologias} resume as tecnologias adotadas em cada camada do \sistema{} \cite{go,sqlc,postgresql,angular}.

\begin{table}[htbp]
  \centering
  \caption{Tecnologias adotadas no \sistema}
  \label{tab:tecnologias}
  \begin{tabularx}{\textwidth}{llX}
    \toprule
    \textbf{Camada} & \textbf{Tecnologia} & \textbf{Função no projeto} \\
    \midrule
    Back-end       & Go 1.27        & API REST e regras de negócio \\
    Acesso a dados & sqlc           & Geração de código Go tipado a partir de consultas SQL \\
    Banco de dados & PostgreSQL 17  & Persistência dos dados \\
    Front-end      & Angular 22     & Interface web \\
    Ambiente       & Docker Compose & Execução opcional do PostgreSQL em contêiner \\
    \bottomrule
  \end{tabularx}
  \fonte{Elaborado pelos autores.}
\end{table}

\section{Segurança e Controle de Acesso}
\orientacao{Descrever as decisões relacionadas à autenticação, autorização, controle de acesso, proteção de dados e demais mecanismos de segurança aplicáveis ao sistema.}

\diagrama{Diagrama de estados}{estados}
